%!TEX root = ../template.tex

\typeout{NT FILE chapter7.tex}%

\chapter{Conclusion and Future Work}
\label{cha:conclusion}

\section{Conclusion}
\label{sec:conclusion}

Conflict-Free Replicated Data Types promise convergence without coordination, but the two established ways of verifying
that promise make different trade-offs. Automated tools built on SMT solvers, VeriFx among them, prove algebraic properties
quickly and with little manual effort, but struggle precisely where a CRDT's correctness depends on recursive definitions
or unbounded structures. Deductive platforms such as Why3 handle exactly this kind of reasoning, but ask the programmer
to write and maintain the proof by hand, in a language with its own syntax and its own demands. Neither approach, used
alone, covers what the other is good at.

This dissertation addressed that gap with a single specification language, compiled automatically to both VeriFx and Why3
from the same source. The language hides the syntactic differences between the two targets behind one set of constructs,
interfaces, axioms, and a handful of directives that let a specification diverge between backends where it needs to, and
the compiler derives an idiomatic implementation for each, an auxiliary and main module pair or a set-based predicate for
Why3, a class extending a shared proof trait for VeriFx, from the same abstract syntax tree.

The evaluation, across ten CRDTs of increasing complexity, showed this design working as intended. Specification effort
dropped for half of the ten, most sharply for the CRDTs built around sets, where a single native collection method absorbs
logic the specification states explicitly. Every one of the ten is proved automatically in Why3, Z3 or Alt-Ergo closing
every obligation on their own except one, \texttt{PNCounter\_CRDT}'s \texttt{compare\_correct}, where a single manual
assertion was enough to split a goal that the solver could not reach on its own into two the solver closed immediately.
VeriFx proves nine of the same ten CRDTs just as automatically, but rejects the tenth, \texttt{RGA}, the one CRDT whose
correctness genuinely depends on reasoning about an unbounded recursive structure, exactly the case the literature already
identifies as beyond the reach of SMT solving alone. Because the same specification also compiles to Why3, \texttt{RGA}
still gets proved, an inductive lemma over the same \texttt{fuel} parameter the compiler strips out for VeriFx is enough
for Why3 to close the proof VeriFx cannot. \texttt{RGA} is, in that sense, the clearest demonstration of why this
dissertation paired the two approaches instead of choosing between them.

\section{Future Work}
\label{sec:future-work}

\texttt{RGA} is the only one of the ten CRDTs whose specification needed the \texttt{fuel} and \texttt{variant} pattern
to give Why3 a termination measure to induct on, and the only one whose correctness depends on reasoning about a recursive
structure of unbounded size in the first place. Extending the evaluation to further recursive or unbounded CRDTs, sequence
CRDTs other than \texttt{RGA} among them, would test whether this same pattern carries over with little change, or whether
it surfaces cases where the language or the compiler would need to grow to support them.
